\begin{recipebox}{Pavlova}{8}{20 min + 1 hr 15 min baking + cooling}
\begin{ingredients}
  \item 4 egg whites, at room temperature
  \item 1 cup caster sugar
  \item 1 tsp white vinegar
  \item 1 tsp vanilla essence
  \item 1 tsp cornflour
  \item[] \textbf{To top}
  \item 300ml cream, whipped
  \item Seasonal fruit: kiwifruit, strawberries, passionfruit or berries
\end{ingredients}

\begin{method}
  \item Preheat oven to 120\degree C bake (not fan). Draw a 20cm circle on a sheet of baking paper and put it upside down on a tray.
  \item Beat the egg whites in a clean, dry bowl until soft peaks form.
  \item Add the sugar a spoonful at a time, beating well after each. Keep beating for 8--10 minutes until very thick and glossy and the sugar has dissolved (rub a little between your fingers --- it shouldn't feel gritty).
  \item Fold in the vinegar, vanilla and cornflour.
  \item Spoon onto the circle and shape into a round with slightly raised sides.
  \item Bake for 1 hour 15 minutes. Turn the oven off and leave the pavlova inside with the door closed until completely cold (ideally several hours or overnight).
  \item Just before serving, top with whipped cream and fruit.
\end{method}

\begin{notes}
Recipes vary on temperature (Edmonds uses 150\degree C for 45 minutes with 6 whites; Chelsea uses 110\degree C). 120\degree C for this 4-white size is a common middle ground.
Don't open the oven door while it bakes or cools --- sudden temperature changes make it crack and collapse.
An un-topped pavlova keeps in an airtight container for 2--3 days. Top only when serving, or it goes soggy.
\end{notes}
\end{recipebox}
